\documentclass[10pt,a4paper]{article}

\usepackage[utf8]{inputenc}
\usepackage[a4paper,margin=0.65in]{geometry}
\usepackage{titlesec}
\usepackage{amsmath,amssymb}
\usepackage{booktabs,tabularx}
\usepackage{xcolor}
\usepackage{hyperref}
\usepackage{enumitem}
\usepackage{fancyhdr}
\usepackage{setspace}

\definecolor{primaryblue}{HTML}{1A3C6E}
\definecolor{accentorange}{HTML}{F97316}
\definecolor{darkgray}{HTML}{1E293B}
\hypersetup{colorlinks=true,linkcolor=primaryblue,urlcolor=primaryblue,citecolor=primaryblue}
\setstretch{1.0}
\setlength{\parindent}{0pt}
\setlength{\parskip}{3pt}
\setlist{nosep,leftmargin=1.2em}
\titleformat{\section}{\large\bfseries\color{primaryblue}}{\thesection}{0.5em}{}
\titleformat{\subsection}{\normalsize\bfseries\color{darkgray}}{\thesubsection}{0.5em}{}
\titlespacing*{\section}{0pt}{7pt}{3pt}
\titlespacing*{\subsection}{0pt}{5pt}{2pt}
\pagestyle{fancy}
\fancyhf{}
\lhead{\textcolor{darkgray}{Aaple Sarkar Portal Redesign}}
\rhead{\textcolor{primaryblue}{HMI Project}}
\cfoot{\thepage}
\setcounter{secnumdepth}{2}

\begin{document}

\begin{center}
{\LARGE\bfseries\color{primaryblue} AAPLE SARKAR 2.0}\\[2pt]
{\large\bfseries Maharashtra Citizen Public Services Portal}\\[4pt]
{\normalsize Redesigning the GUI to Minimize Screen Complexity and Cognitive Entropy}\\[6pt]
\rule{0.9\textwidth}{1pt}\\[5pt]
{\large\bfseries Human-Machine Interaction (HMI)}\\
{\normalsize Project / Practical 8 Report}\\[5pt]
{\small \textbf{Approach:} Goal-Directed Design (GDD) and Tullis Complexity Reduction}\\[6pt]
\textbf{Submitted by}\\
Sanika Lobo (10613 -- Team Leader) \quad Meet Mangaonkar (10616)\\
Yash Masaye (10617) \quad Jeet Chavan (10584)\\[4pt]
\textbf{Department of Computer Engineering}\\
Fr. Concei\c{c}\~ao Rodrigues College of Engineering\\
Academic Year 2026 \quad | \quad Semester V
\rule{0.9\textwidth}{0.6pt}
\end{center}

\section*{Abstract}
The Aaple Sarkar portal (\url{https://aaplesarkar.mahaonline.gov.in}) provides Maharashtra citizens access to statutory services such as Income, Caste, Domicile and Solvency certificates under the Right to Public Services (RTS) Act, 2015. Its interface presents usability issues including multiple navigation bars, oversized banners, early CAPTCHA placement, moving notices and a dense collection of links. This project evaluates the interface using screen-complexity concepts and proposes a six-screen redesign based on Goal-Directed Design, progressive disclosure and clear feedback. The reported comparison reduces alignment lines from 73 to 14 (80.8\%), visual elements from 148 to 34 (77.0\%), and estimated visual search time from 3.42 to 0.94 seconds (72.5\%).

\section{Introduction and Contextual Inquiry}
Government-to-Citizen (G2C) portals must serve users with different levels of digital literacy and devices. Although Aaple Sarkar supports access to important government services, the interface can make service discovery and application submission unnecessarily difficult.

\textbf{Observed usability issues}
\begin{enumerate}
\item \textbf{Navigation redundancy:} three stacked navigation bars, repeated logos and confusing menus.
\item \textbf{Poor use of screen space:} large banners and photographs push important services below the fold.
\item \textbf{Premature CAPTCHA:} security verification appears before a user chooses a service or begins authentication.
\item \textbf{Visual competition:} uneven sidebars and auto-scrolling notices compete for attention and reduce user control.
\item \textbf{Weak grouping:} more than 48 links appear without clear categories, increasing visual search effort.
\item \textbf{Fragmented account recovery:} separate recovery options and legacy notices add clutter.
\end{enumerate}

\section{Screen Complexity Models}
\textbf{Information entropy.} A visual-disorder measure can be expressed as
\[
C=-N\sum_{i=1}^{k}p_i\log_2(p_i), \qquad p_i=\frac{n_i}{N}
\]
where \(N\) is the number of interface elements, \(k\) is the number of alignment coordinates, and \(n_i\) is the number of elements sharing coordinate \(i\). The probabilities \(p_i\) describe how elements are distributed across alignments.

\textbf{Tullis alignment complexity.} Alignment complexity is represented as
\[
C_{\mathrm{alignment}}=C_h+C_v
\]
where \(C_h\) and \(C_v\) are the numbers of distinct horizontal and vertical alignment starting points. A simplified visual-search estimate is
\[
T\approx k_1+k_2\log_2(N)+k_3(C_h+C_v)
\]
where the constants depend on the measurement method. Fewer elements and consistent alignment are expected to reduce the effort required to locate controls.

\section{Before-and-After Evaluation}
The following figures are the project's reported estimates for the existing and redesigned interfaces.

\begin{center}
\small
\begin{tabularx}{\textwidth}{|>{\raggedright\arraybackslash}X|c|c|c|}
\hline
\textbf{Metric} & \textbf{Existing} & \textbf{Redesign} & \textbf{Change}\\
\hline
Total visual elements & 148 & 34 & 77.0\% fewer\\
\hline
Vertical alignment lines ($C_v$) & 31 & 6 & 80.6\% fewer\\
\hline
Horizontal alignment lines ($C_h$) & 42 & 8 & 80.9\% fewer\\
\hline
Total alignment complexity & 73 & 14 & 80.8\% lower\\
\hline
Estimated visual search time & 3.42 s & 0.94 s & 72.5\% lower\\
\hline
Navigation tiers & 3 & 1 sticky header & 66.7\% fewer\\
\hline
Homepage choices & 48 links & 5 actions + search & 87.5\% fewer links\\
\hline
\end{tabularx}
\end{center}

The proposed layout uses a consistent grid, more whitespace and grouped service cards. Decorative photographs, animated tickers, duplicated seals and legacy badges are removed. These changes are intended to create a clearer scan path and reduce competing visual targets. The search-time values are model-based estimates, not results of a controlled user study.

\section{Redesigned Interface: Six Task-Focused Screens}
The redesign separates the portal into screens organized around citizens' main tasks.

\textbf{1. Homepage}
\begin{itemize}
\item One sticky header with the Aaple Sarkar emblem, Home, Services, Track, Help \& FAQ, Contact and Citizen Login.
\item Deep-blue hero area with a prominent service-search field, saffron action button and popular-service shortcuts.
\item Three quick-action cards: Apply for Service, Track Application and Download Certificate.
\item Department filters and a responsive service-card grid; a compact strip highlights Aadhaar integration, DigiLocker, PayGov and digital signatures.
\end{itemize}

\textbf{2. All Services Catalog}\\
A searchable, filterable catalog groups the 13 proposed essential services: Income, Caste, Age/Nationality/Domicile, Senior Citizen, Solvency, Non-Creamy Layer, Landless, General Affidavit, Certified Copy (7/12), Agriculturist, Cultural Programme Permission, Small Land Holder and Residence in Hilly Area certificates. Each card provides an application action and, where applicable, an expected service timeline.

\textbf{3. Login and Registration}\\
A segmented choice separates Citizen Login from VLE (Kendra) Login. CAPTCHA appears only during authentication and includes a refresh control. A single “Forgot credentials?” link opens OTP-based recovery. Registration groups district, Aadhaar and mobile verification into a clear sequence.

\textbf{4. Application Tracking}\\
A centered tracking-ID field allows lookup. The result card presents relevant application details and a status label such as Approved, Under Review, Pending or Rejected. A four-stage progress indicator shows Submitted, Departmental Review, Decision Made and Certificate Ready. Available certificates and acknowledgements can be downloaded.

\textbf{5. Three-Step Application Wizard}
\begin{enumerate}
\item \textbf{Personal details:} name, date of birth, gender, address, district and taluka.
\item \textbf{Document upload:} clear upload areas, file-size guidance and a proposed DigiLocker retrieval option.
\item \textbf{Review and submit:} summary of entered details, applicable fee, declaration and confirmation receipt with a tracking ID.
\end{enumerate}
A prerequisite checklist is shown before users begin, reducing the chance of discovering missing documents midway through the process.

\textbf{6. Help and Support}\\
Searchable FAQs are grouped into General, Login, Process, Documents and Downloads. A support card provides the verified helpline, official email and service hours. Contact information should be confirmed against the current official portal before release.

\section{HMI Principles Applied}
\begin{itemize}
\item \textbf{Hick--Hyman Law:} reducing the number of competing choices supports faster decisions.
\item \textbf{Fitts's Law:} sufficiently large buttons and prominent primary actions make controls easier to target.
\item \textbf{Miller's Law and chunking:} grouping information and dividing long forms into three stages makes tasks easier to manage.
\item \textbf{Norman's feedback and visibility:} status labels, progress indicators, loading feedback and confirmation messages make system state clear.
\item \textbf{Progressive disclosure:} CAPTCHA and detailed instructions appear when relevant rather than cluttering the homepage.
\item \textbf{Accessibility:} the proposed interface includes readable contrast, font-size controls, English/Marathi language selection and clear labels. Accessibility should be validated against WCAG 2.1 AA.
\end{itemize}

\section{Conclusion}
The Aaple Sarkar redesign demonstrates how a task-focused structure can make a government-services interface easier to scan and navigate. The proposal consolidates navigation, removes unnecessary visual elements, groups services, and guides applicants through tracking and form submission. The reported comparison shows an 80.8\% reduction in alignment complexity, 77.0\% fewer visual elements and a 72.5\% reduction in estimated search time. These values should be treated as design estimates until verified through usability testing with representative citizens. The six-screen structure aims to preserve essential functionality while providing a more consistent, accessible and citizen-centered experience.

\end{document}